\label{sec:appendix-pto-model}
To implement the multiport model of \Cref{sec:dynamics}, we must specify its
conventions, PTO kinematics, PTO dynamics, power relations, and ultimate solution procedure.
This section provides a general structure in effort-flow form applicable across multiple energy domains for any number of degrees of freedom, as well as specific equations for a common PTO structure consisting of a mechanical drivetrain and a permanent magnet generator.

%%%%%%%%%%%%%%%
\subsubsection{Multiport Circuit Conventions}
\label{sec:appendix-multiport}

The intrinsic, powertrain, and Th\'evenin-equivalent representations used throughout \Cref{sec:dynamics} are all instances of multiport circuit models,
with effort variables $\gls{sym-vec-e}$ and flow variables $\gls{sym-vec-q}$ at each port.
Following \citet{reveyrand_multiport_2018}, three matrix forms describe the relationships between port variables:
impedance ($\gls{sym-mathbf-Z}$), cascade ($\gls{sym-mathbf-a}$, also called ABCD), and hybrid ($\gls{sym-mathbf-h}$):
\begin{equation}\label{eq:multiport-matrix-definitions}
    \vec{e} = \gls{sym-mathbf-Z}\,\gls{sym-vec-q},
\qquad
    \begin{bmatrix}
        \vec{e}_2 \\ \vec{q}_2
    \end{bmatrix}
    = [\gls{sym-mathbf-a}]_{2\leftarrow 1}
    \begin{bmatrix}
        \vec{e}_1 \\ \vec{q}_1
    \end{bmatrix},
\qquad
    \begin{bmatrix}
        \vec{e}_1 \\ \vec{q}_2
    \end{bmatrix}
    = [\gls{sym-mathbf-h}]_{1,2}
    \begin{bmatrix}
        \vec{q}_1 \\ \vec{e}_2
    \end{bmatrix}.
\end{equation}
These equations apply to both time-domain and phasor representations of each port.
The three forms differ in how the ports are grouped:
\begin{itemize}
    \item \textbf{Impedance} ($\gls{sym-mathbf-Z}$)
        groups all ports together, describing the effort required to produce a unit of flow.
        The intrinsic dynamics matrix $\gls{sym-mathbf-Z-i}$ and the PTO impedance
        $\gls{sym-mathbf-Z-p}$ in \Cref{eq:eom-freq-domain,eq:intrinsic-impedance} both use this form.
    \item \textbf{Cascade} ($\gls{sym-mathbf-a}$)
        divides the ports into two groups and describes how effort and flow transmit from one group to the other.
        The arrow subscript indicates direction:
        $[\gls{sym-mathbf-a}]_{2\leftarrow 1}$ maps port group 1 to port group 2.
        Reversing direction inverts the matrix,
        \begin{equation}\label{eq:cascade-inverse}
            [\gls{sym-mathbf-a}]_{1\leftarrow 2} = [\gls{sym-mathbf-a}]_{2\leftarrow 1}^{-1},
        \end{equation}
        and sequential multiplication represents serial composition of two-port elements.
        This form is used for the PTO dynamics (drivetrain followed by generator).
    \item \textbf{Hybrid} ($\gls{sym-mathbf-h}$)
        mixes the two conventions, expressing one effort and one flow at the two port groups as functions of the complementary variables.
        This form is used for the PTO kinematics, where it naturally accommodates underactuation.
\end{itemize}

Conversion between the three forms is accomplished using the formulas proposed by \citet{reveyrand_multiport_2018}.
Some port relationships are physically well-defined in one convention but degenerate in another (for example, an ideal current source has no impedance representation but a well-defined hybrid representation) so the convention is chosen to match the physics of each subsystem.

% Verify against \citet{coe_co-design_2025} before submission and edit
% if the convention differs.
The sign convention adopted here treats both flows as positive when entering the multiport,
matching the suggestion of \citet{coe_co-design_2025}.

%%%%%%%%%%%%%%%
\subsubsection{PTO Kinematics and Underactuation}
We next turn our attention to the PTO kinematics and dynamics, to specify the generic impedance matrix $\gls{sym-mathbf-Z-p}$ in \Cref{eq:eom-freq-domain}.
A PTO model must capture the dynamics of any elements that transmit, transduce, or condition power between the rigid-body mechanical energy and the end energy product, in this case electricity.
Relevant dynamics include the inertia, stiffness, static friction, mechanical advantage, losses, and electromagnetic coupling of flywheels, shafts, springs, gears, hydraulic circuits, generators, and power electronics.
\citet{penalba_review_2016} review WEC PTO models and \citet{coe_co-design_2025} synthesize these into a multiport impedance matching network framework to facilitate control co-design.
\citet{coe_co-design_2025} present a PTO model in cascade matrix form for a single degree of freedom system.
The PTO kinematics describe the relationship between WEC body displacements $\gls{sym-vec-hat-xi}$ to PTO displacements $\gls{sym-vec-hat-X-PTO}$,
determined by linkages or similar mechanisms and not including subsequent PTO elements such as a mechanical drivetrain or generator.
In general form, the PTO kinematics can be expressed as:
\begin{equation}\label{eq:pto-kinematics-motion}
\gls{sym-vec-hat-X-PTO} = \gls{sym-mathbf-T-kin}~ \gls{sym-vec-hat-xi}
\end{equation}
with kinematic transformation matrix $\gls{sym-mathbf-T-kin}$ that may depend on the body displacements $\gls{sym-vec-hat-xi}$ if the linkages are nonlinear,
but hereafter is assumed constant for simplicity.
This assumption means that \Cref{eq:pto-kinematics-motion} can be differentiated trivially to yield the same matrix relationship for velocities.
Assuming power conservation across the linkages,
the kinematics describe the relationship between not only motions but also forces:
\begin{equation}\label{eq:pto-kinematics-force}
\gls{sym-vec-hat-F-p} = [\gls{sym-mathbf-T-kin}]^\top  \gls{sym-vec-hat-F-PTO}
\end{equation}
To avoid confusion, subscript $PTO$ is used for forces in the PTO degrees of freedom (post-kinematics),
while subscript $p$ denotes forces in the hydrodynamic body degrees of freedom (pre-kinematics).
\Cref{eq:pto-kinematics-motion,eq:pto-kinematics-force} imply rigid lossless massless linkages,
and any linkage deflection, dissipation, and inertia would be included in the PTO dynamics rather than kinematics.

Multibody WECs are frequently underactuated,
meaning that the PTO has fewer controlled degrees of freedom than the hydrodynamic bodies,
and $\gls{sym-mathbf-T-kin}$ has more columns than rows.
Underactuated systems receive considerable attention in robotics \citep{tedrake_underactuated_2024}
but have only recently been explored in the context of wave energy conversion \citep{faedo_principle_2022}.
Cascade matrices require an identical number of input and output ports
and thus cannot be used to describe the unbalanced ports that underactuated kinematics create
\citep{reveyrand_multiport_2018}.
Non-square impedance matrices may be used to capture unbalanced port relations in general,
but the impedance is undefined for the specific relations of
\Cref{eq:pto-kinematics-motion,eq:pto-kinematics-force}.
Instead, the kinematics can be packaged into a square hybrid matrix $\gls{sym-mathbf-h-kin}$
which accurately represents the underactuated system:
\begin{equation}\label{eq:h-matrix-kinematics}
    \begin{bmatrix}
        \gls{sym-vec-hat-F-p} \\
        \gls{sym-vec-hat-dot-X-PTO}
    \end{bmatrix}
    = \underbrace{
        \begin{bmatrix}
            \mathbf{0}             & \gls{sym-mathbf-T-kin}^{\top} \\
            \gls{sym-mathbf-T-kin} & \mathbf{0}
        \end{bmatrix}}
        _{\gls{sym-mathbf-h-kin}}
    \begin{bmatrix}
        \gls{sym-vec-hat-dot-xi} \\ \gls{sym-vec-hat-F-PTO}
    \end{bmatrix}
\end{equation}

Here, we model two hydrodynamic degrees of freedom (heave of the float and spar respectively) and a single PTO degree of freedom corresponding to the relative heave motion between the float and spar.
The kinematics matrix is therefore:
\begin{equation}
    \label{eq:T-kin-specific}
    \gls{sym-mathbf-T-kin} = \begin{bmatrix}1 & -1\end{bmatrix}
\end{equation}
Thus, even absent any mechanism with ``kinematics'' in the traditional sense,
the kinematics matrix is filled with zeroes, ones, and negative ones to represent which degrees of freedom are controlled.

In the circuit representation with two hydrodynamic and one PTO degree of freedom,
the hybrid matrix \Cref{eq:h-matrix-kinematics} can be represented as a three-port element,
as shown in \Cref{fig:multiport-circuit-kinematics}.
\begin{figure}[htbp]
    \centering
    \includegraphics[\FigureOptions{aor:figs/from-matlab/circuit_kinematics.pdf}]{figs/from-matlab/circuit_kinematics.pdf}
    \caption{Multiport circuit with powertrain kinematics represented as a hybrid matrix, followed by the Th\'{e}venin equivalent of the intrinsic dynamics and kinematics}
    \label{fig:multiport-circuit-kinematics}
\end{figure}

%%%%%%%%%%%%%%%
\subsubsection{Combining Intrinsic Dynamics with PTO Kinematics}
Taking the Th\'{e}venin equivalent of the intrinsic dynamics and kinematics at the PTO port yields a scalar equivalent circuit for each PTO degree of freedom,
with force sources $\gls{sym-vec-hat-F-th}$ and mechanical impedances $\gls{sym-vec-Z-i-th}$,
shown in the bottom portion of \Cref{fig:multiport-circuit-kinematics} and computed with the following expressions:
\begin{equation}
\begin{aligned}
    \gls{sym-vec-hat-F-th} &= (\gls{sym-mathbf-T-kin}\gls{sym-mathbf-T-kin}^{\top})^{-1}
                               \gls{sym-mathbf-T-kin} \gls{sym-mathbf-Z-i}^{-1} \gls{sym-vec-gamma} \gls{sym-hat-zeta}
    = \frac{1}{2} [1,~-1] \gls{sym-mathbf-Z-i}^{-1} \gls{sym-vec-gamma} \gls{sym-hat-zeta} \\
    \gls{sym-vec-Z-i-th} &= (-\gls{sym-mathbf-T-kin}\gls{sym-mathbf-Z-i}^{-1} \gls{sym-mathbf-T-kin}^{\top})^{-1}
                          = -\frac{\det(\gls{sym-mathbf-Z-i})}{\Sigma\gls{sym-mathbf-Z-i}}
                          = -\frac{Z_{i,f}Z_{i,s}-Z_{i,c}^2}{Z_{i,f}+Z_{i,s}+2Z_{i,c}}
\end{aligned}
\end{equation}
where the vector expressions for the general case are shown first,
followed by the scalar result for the $\gls{sym-mathbf-T-kin}$ in \Cref{eq:T-kin-specific} with 1 controlled degree of freedom.

Although, as previously mentioned, a true cascade matrix must have the same number of input and output ports, it is possible to formulate a cascade-like matrix that can transform from the PTO port to the body ports.
However, it cannot transform in the other direction due to underactuation.

Even in the favorable direction, this matrix cannot transform to the body ports directly.
Due to the wave excitation, body velocities can be nonzero even when the PTO force and velocity are both zero.
Hence, there is no way to express $\gls{sym-hat-dot-xi}$ as a weighted sum of $\gls{sym-hat-F-PTO}$ and $\gls{sym-dot-hat-X-PTO}$.
$\gls{sym-hat-dot-xi}$ could be expressed as a function of the PTO variables by substituting
$\gls{sym-mathbf-Z-p}=\gls{sym-mathbf-T-kin}^{\top} (\gls{sym-vec-hat-F-PTO}/\gls{sym-vec-hat-dot-X-PTO}) \gls{sym-mathbf-T-kin}$
into \Cref{eq:eom-freq-domain}, but this yields a nonlinear expression of limited utility.
Instead, we can use cascade-like matrices to solve for a modified variable $\gls{sym-vec-dot-xi-Delta}$ defined as the change in the body velocities caused by the PTO,
in other words the difference between the actual body velocities $\gls{sym-vec-dot-xi}$
and their values if the PTO were disconnected with $\gls{sym-F-PTO}=0$, $\gls{sym-vec-dot-xi-0}$:
\begin{equation}\label{eq:modified-body-velocity}
    \gls{sym-vec-dot-xi-Delta} = \gls{sym-vec-dot-xi} - \gls{sym-vec-dot-xi-0}
    \quad \textrm{ where } \quad
    \gls{sym-vec-dot-xi-0} = \gls{sym-mathbf-Z-i}^{-1}~\gls{sym-vec-gamma}~\frac{\gls{sym-H}}{2}
\end{equation}
The cascade-like matrix from the PTO port $\{\gls{sym-F-PTO}, \gls{sym-dot-X-PTO}\}$
to the modified body port $\{\gls{sym-F-p}, \gls{sym-dot-xi-Delta}\}$ is then:
\begin{equation}\label{eq:body-transmission-matrix}
    [\gls{sym-mathbf-a}]_{\gls{sym-F-p}\gls{sym-xi-Delta}\leftarrow \gls{sym-F}\gls{sym-dot-X}} =
    \begin{bmatrix}
        \gls{sym-mathbf-I}^{2\times2} \\
        -\gls{sym-mathbf-Z-i}^{-1}
    \end{bmatrix}
    \gls{sym-mathbf-T-kin}^{\top}
    \begin{bmatrix}
        1 & 0
    \end{bmatrix}
\end{equation}
noting that this matrix is non-square ($\in \gls{sym-mathbb-C}^{4\times2}$) since there are more hydrodynamic than PTO degrees of freedom,
and that the matrix label drops the $PTO$ subscripts for compactness.
\Cref{fig:multiport-circuit-kinematics} represents this in circuit form,
noting that the excitation source must be duplicated to create $\gls{sym-vec-hat-dot-xi-0}$.

%%%%%%%%%%%%%%%
\subsubsection{PTO Dynamics}
\label{sec:appendix-pto-dynamics}
In addition to the PTO port already discussed, we define two additional ports of relevance:
the generator mechanical port (torque $\gls{sym-tau}$ and rotation speed $\gls{sym-Omega}$)
and the generator electrical port (voltage $\gls{sym-V}$ and current $\gls{sym-I}$).
Assuming a drivetrain consisting of a linear to rotational gear ratio with equivalent radius $\gls{sym-R-gear}$ followed by a driveshaft with impedance $\gls{sym-Z-shaft}$,
the PTO mechanical dynamics are expressed with cascade matrix
$[\gls{sym-mathbf-a-F-dot-X-leftarrow-tau-Omega}]$:
\begin{equation}\label{eq:mech-transmission-matrix}
\begin{bmatrix}
\gls{sym-F-PTO} \\ \gls{sym-dot-X-PTO}
\end{bmatrix}
=
\underbrace{\begin{bmatrix}
\gls{sym-Z-shaft}/\gls{sym-R-gear} & 1/\gls{sym-R-gear} \\
\gls{sym-R-gear} & 0
\end{bmatrix}}
_{[\gls{sym-mathbf-a-F-dot-X-leftarrow-tau-Omega}]}
\begin{bmatrix}
 \gls{sym-tau}\\ \gls{sym-Omega}
\end{bmatrix}
\end{equation}

Meanwhile, a synchronous permanent magnet generator modeled as an ideal gyrator with torque constant $\gls{sym-k-tau}$ plus a winding impedance $\gls{sym-Z-w}$ has cascade matrix $[\gls{sym-mathbf-a}]_{\gls{sym-tau} \gls{sym-Omega} \leftarrow \gls{sym-V} \gls{sym-I}}$:
\begin{equation}\label{eq:elec-transmission-matrix}
    \begin{bmatrix}
        \gls{sym-tau} \\ \gls{sym-Omega}
    \end{bmatrix}
    =
    \underbrace{\begin{bmatrix}
    0 & \gls{sym-k-tau} \\
    1/\gls{sym-k-tau} & \gls{sym-Z-w}/\gls{sym-k-tau}
    \end{bmatrix}}
    _{[\gls{sym-mathbf-a-tau-Omega-leftarrow-V-I}]}
    \begin{bmatrix}
        \gls{sym-V} \\ \gls{sym-I}
    \end{bmatrix}
\end{equation}
with Park-transformed quadrature electrical voltage and current $\gls{sym-V}$ and $\gls{sym-I}$.

The system operating point is determined by a controller that applies some electrical load to the generator.
The controller can be formulated as an impedance at either the generator mechanical or electrical port.
The former is more convenient for hardware implementation, because modern drives typically expect torque commands,
while the latter is more convenient for electrical power calculations and optimal control.
%Here we formulate both as linear impedances with damping and stiffness terms, where
The mechanical control impedance $\gls{sym-Z-u}$ and electrical load impedance $\gls{sym-Z-l}$ are defined as follows:
\begin{equation}\label{eq:controller-impedances}
    \gls{sym-Z-l} = \frac{\gls{sym-V}}{\gls{sym-I}} = \gls{sym-B-l} + \frac{\gls{sym-K-l}}{i\gls{sym-omega}}
    ,  \qquad
    \gls{sym-Z-u} = \frac{\gls{sym-tau}}{\gls{sym-Omega}}% = B_u + \frac{K_u}{s}
\end{equation}
Combining \Cref{eq:controller-impedances,eq:elec-transmission-matrix} gives the following expression for $\gls{sym-Z-u}$ as a function of $\gls{sym-Z-l}$:
\begin{equation}\label{eq:control-impedance}
    \gls{sym-Z-u} =
    \frac{
    \begin{bmatrix}
    1 & 0
    \end{bmatrix}
    [\gls{sym-mathbf-a-tau-Omega-leftarrow-V-I}]
    \begin{bmatrix}
    \gls{sym-Z-l} \\ 1
    \end{bmatrix}
}{
    \begin{bmatrix}
    0 & 1
    \end{bmatrix}
    [\gls{sym-mathbf-a-tau-Omega-leftarrow-V-I}]
    \begin{bmatrix}
    \gls{sym-Z-l} \\ 1
    \end{bmatrix}
}
\end{equation}
This equation allows the $\gls{sym-Z-l}$ identified with optimal control to be implemented as a corresponding $\gls{sym-Z-u}$ in hardware.
%Incorporating \Cref{eq:mech-transmission-matrix} yields the Th\'{e}venin equivalent scalar PTO impedance at the hydrodynamic body port:
% \begin{equation}\label{eq:pto-impedance}
% Z_{PTO,th} =
% \frac{
%     \begin{bmatrix}
%     1 & 0
%     \end{bmatrix}
%     [\mathbf{a}]_{F\dot{X}\leftarrow \tau \Omega}
%     [\mathbf{a}]_{\tau \Omega \leftarrow V I}
%     \begin{bmatrix}
%     Z_l \\ 1
%     \end{bmatrix}
% }{
%     \begin{bmatrix}
%     0 & 1
%     \end{bmatrix}
%     [\mathbf{a}]_{F\dot{X}\leftarrow \tau \Omega} [\mathbf{a}]_{\tau \Omega \leftarrow V I}
%     \begin{bmatrix}
%     Z_l \\ 1
%     \end{bmatrix}
% }
% \end{equation}
% which is useful for calculating the mechanical response of each body.

The PTO dynamics can be represented as a multiport circuit, as shown in \Cref{fig:multiport-circuit-pto}.
\begin{figure}[htbp]
\centering
    \includegraphics[\FigureOptions{aor:figs/from-matlab/circuit_pto.pdf}]{figs/from-matlab/circuit_pto.pdf}
\caption{PTO dynamics represented as a multiport circuit and electrical Th\'{e}venin equivalent}
\label{fig:multiport-circuit-pto}
\end{figure}

The intrinsic dynamics, PTO kinematics, and PTO dynamics can be combined into a Th\'{e}venin equivalent source voltage and impedance $\gls{sym-hat-V-s-th}$ and $\gls{sym-Z-s-th}$ as follows:
\begin{equation}\label{eq:thevenin-electrical}
    \gls{sym-Z-s-th} = \frac{
        \begin{bmatrix}1 & 0\end{bmatrix}
        \gls{sym-mathbf-a-VI-leftarrow-F-dot-X}
        \begin{bmatrix}-\gls{sym-Z-i-th} \\ 1\end{bmatrix}
    }
    {
        \begin{bmatrix}0 & 1\end{bmatrix}
        \gls{sym-mathbf-a-VI-leftarrow-F-dot-X}
        \begin{bmatrix}-\gls{sym-Z-i-th} \\ 1\end{bmatrix}
    } = \frac{
                \gls{sym-B-a} - \gls{sym-A-a} \gls{sym-Z-i-th}
            }
            {
                \gls{sym-D-a}-\gls{sym-C-a} \gls{sym-Z-i-th}
            }
    \qquad
    \gls{sym-hat-V-s-th} =
        \gls{sym-hat-F-th}
        (\gls{sym-A-a} - \gls{sym-C-a} \gls{sym-Z-s-th})
\end{equation}
where $\gls{sym-A-a},\gls{sym-B-a},\gls{sym-C-a},\gls{sym-D-a}$
are the elements of
$\gls{sym-mathbf-a}_{\gls{sym-V} \gls{sym-I} \leftarrow F\gls{sym-dot-X}}$.
This is shown in the bottom portion of \Cref{fig:multiport-circuit-pto}.
The load voltage and current can then be found as:
\begin{equation}\label{eq:voltage-current}
    \gls{sym-hat-V} = \gls{sym-hat-V-s-th}\frac{
                                                    \gls{sym-Z-l}
                                                }
                                                {
                                                    \gls{sym-Z-s-th} + \gls{sym-Z-l}
                                                },
            \qquad \gls{sym-hat-I} = \frac{
                                                \gls{sym-hat-V-s-th}
                                            }
                                            {
                                                \gls{sym-Z-s-th}+\gls{sym-Z-l}
                                            }
\end{equation}

%%%%%%%%%%%%%%%

\subsubsection{Time-Domain Expansion of Instantaneous Power}
\label{sec:appendix-power-time}

Equation~\eqref{eq:power-PQS} can be used to rewrite the instantaneous power $\gls{sym-p}(\gls{sym-t})$ at any port as \citep{saadat_power_1999}
\begin{equation}\label{eq:power-time-PQS}
    \gls{sym-p}(\gls{sym-t}) =
    \gls{sym-P} + |\gls{sym-S}|
    \cos\!\left(
                2(\gls{sym-omega} \gls{sym-t} + \angle\hat{e})
                - \tan^{-1}\!\left(
                                    \frac{\gls{sym-Q}} {\gls{sym-P}}
                            \right)
        \right),
\end{equation}
a DC offset of $\gls{sym-P}$ plus an oscillation at $2\gls{sym-omega}$ with amplitude $|\gls{sym-S}|$.
The average and peak powers in \Cref{eq:power-avg-peak} follow directly:
$\gls{sym-p-avg} = \gls{sym-P}$ since the oscillation averages to zero,
and $\gls{sym-p-pk} = \gls{sym-P} + |\gls{sym-S}|$ at the cosine maximum.
The minimum $\gls{sym-p-min} = \gls{sym-P} - |\gls{sym-S}|$ is negative when $|\gls{sym-S}|>\gls{sym-P}$,
corresponding to instants when energy flows from the load back into the source under reactive control.

%%%%%%%%%%%%%%%
\subsubsection{Electrical-Port Power in Optimal-Control Variables}
\label{sec:appendix-power-elec}

Expressing the average and peak powers at the electrical port in terms of the constrained-control parameters introduced in \Cref{sec:optimal-control}
(control damping $\gls{sym-B-l}$, control stiffness $\gls{sym-K-l}$, and current phasor $\gls{sym-hat-I}$),
and applying the constant electrical efficiency $\gls{sym-eta}$,
\begin{equation}\label{eq:power-elec}
    \begin{bmatrix}
        \gls{sym-p-avg-elec} \\
        \gls{sym-p-pk-elec}
    \end{bmatrix}
    =
    \gls{sym-eta}
    \begin{bmatrix}
        \gls{sym-p-avg-VI} \\
        \gls{sym-p-pk-VI}
    \end{bmatrix}
    =
    \gls{sym-eta} \tfrac{1}{2} \gls{sym-B-l} |\gls{sym-hat-I}|^2
    \begin{bmatrix}
        1 \\
        1 + \sqrt{1 + \left(
                        \tfrac{\gls{sym-K-l}} {\gls{sym-B-l}\gls{sym-omega}}
                    \right)^2}
    \end{bmatrix}.
\end{equation}
The peak-power form is used as the apparent-power constraint in the QCQP of \Cref{sec:optimal-control}.

%%%%%%%%%%%%%%%
% Additionally, the drivetrain mechanical cascade matrix $[\gls{sym-a}]_{\gls{sym-m}}$ is set to the identity matrix, which represents a static effective gear radius $\gls{sym-R}=1$:
% \begin{equation}
% [a]_m = \begin{bmatrix}
% 1/R &~ 0 \\
% 0 &~ R
% \end{bmatrix} = \begin{bmatrix}
% 1 &~ 0 \\
% 0 &~ 1
% \end{bmatrix}
% \end{equation}
% With this $[\gls{sym-a}]_{\gls{sym-m}}$, the control damping and stiffness $\gls{sym-B-u}$ and $\gls{sym-K-u}$ equal the PTO damping and stiffness $\gls{sym-B-p}$ and $\gls{sym-K-p}$.
% This will simplify the impedance matching procedure described below, requiring matching at only one interface.

%The time-domain mechanical power is $\gls{sym-P}_{\text{mech}}(\gls{sym-t}) = \gls{sym-tau}(\gls{sym-t}) \gls{sym-Omega}(\gls{sym-t}) = (\gls{sym-B-u} \gls{sym-Omega}(\gls{sym-t}) + \gls{sym-K-u} \gls{sym-theta}(\gls{sym-t}) ) \gls{sym-Omega}(\gls{sym-t}) $, where $\gls{sym-theta}(\gls{sym-t})$ is the generator rotational position. Assuming sinusoidal waveforms, the time-average and peak mechanical power in a given sea state are then
%%%%%%%%%%%%%%%

%%%%%%%%%%%%%%%%%%%%%%%%%%%%%%%%%%%%%%%%%
\subsubsection{Linear Solution Procedure}
\label{sec:appendix-dynamics-solution}
\Cref{sec:eom,sec:pto-dynamics-overview} introduce linear models for the rigid body dynamics, PTO kinematics, and PTO dynamics
using impedance, hybrid, and cascade matrices respectively.
Here we unify the formulation and present a clear system-level solution procedure for all port variables:
\begin{enumerate}
    \item Use \Cref{eq:thevenin-electrical} to assemble the Th\'{e}venin source voltage and impedance
        $\gls{sym-hat-V-s-th}$ and $\gls{sym-Z-s-th}$
        as seen from the port on which power should be maximized
        (in this case, the electrical generator port
        $\{\gls{sym-V},\gls{sym-I}\}$).
    \item Set the load impedance $\gls{sym-Z-l}$.
        Any load impedance can be simulated,
        but typically it is selected as a function of $\gls{sym-Z-s-th}$ to maximize the power at this port,
        possibly subject to constraints on various ports.
        \Cref{sec:optimal-control} frames the load impedance selection as a constrained optimization problem.
    \item From  $\gls{sym-hat-V-s-th}$, $\gls{sym-Z-s-th}$, and $\gls{sym-Z-l}$,
        calculate the $\{\gls{sym-V},\gls{sym-I}\}$ port variables from \Cref{eq:voltage-current}.
    \item Use cascade matrices \Cref{eq:mech-transmission-matrix,eq:elec-transmission-matrix,eq:body-transmission-matrix}
        to propagate the variables back toward the rigid body,
        yielding the $\{\gls{sym-tau},\gls{sym-Omega}\}$, $\{\gls{sym-F-PTO},\gls{sym-dot-X-PTO}\}$,
        and $\{\gls{sym-F-p}, \gls{sym-dot-xi-Delta}\}$ port variables.
    \item Transform from $\gls{sym-dot-xi-Delta}$ to $\gls{sym-dot-xi}$ via \Cref{eq:modified-body-velocity}.
\end{enumerate}
Once all the port variables have been calculated,
power at each port can be calculated according to \Cref{sec:power}.
This serves as the solution to a linear subproblem,
which is then re-solved using updated quasi-linear coefficients until drag and force saturation nonlinearities have converged.
